\uuid{cgwl}
\titre{Matrice/Somme/Produit/Puissance}
\auteur{hanani ; mzari}
\organisation{exo7}
\licence{CC BY-NC-SA 4.0}
\datecreate{2020-01-20}
\niveau{L1}
\module{Algèbre}
\chapitre{Matrice}
\sousChapitre{Propriétés élémentaires, généralités}
\difficulte{2}
\theme{Calcul matriciel, Matrice, Somme, Produit, Puissance}
\exo7qcmid{lille-2.matrices.008}

\contenu{
\begin{question}
\qclassification{108.01}
\qtags{motcle=Matrice/Somme/Produit/Puissance}
Soit $A=\left(\begin{array}{rc}
1&1\\1&1\end{array}\right) $ et $I=  \left(\begin{array}{rc}
1&0\\0&1\end{array}\right) $, la matrice identité.
Quelles sont les assertions vraies ?
\begin{answers}
\good{$A^2=2A$}.
\bad{$A^n=2^nA$, pour tout entier $n \ge 1$}.
\good{$(A-I)^{2n}= I$, pour tout entier $n \ge 1$}.
\bad{$(A-I)^{2n+1}= A+I$, pour tout entier $n \ge 1$}.
\end{answers}
\begin{explanations} On vérifie que pour tout entier $n\ge 1$,
$A^n=2^{n-1}A\, , \, (A-I)^{2n}= I$ et $(A-I)^{2n+1}= A-I$.
\end{explanations}
\end{question}
}
